% Section 8 -- Conclusions. PR: longer than the abstract; strengths, weaknesses, benefit to others, future work.
\section{Conclusions}\label{sec:conclusions}

When each site contributes a single diagnosis, subject-level cross-validation rewards models that score by site, and without an assumption that the data cannot test, diagnosis and site cannot be separated at all. We derived three rules and tested each on three cohorts. Holding out whole sites removed an optimism of $0.37$ to $0.46$ in AUC. When bridging subjects were added, what mattered was that they crossed the class boundary, and $L$, computed from the site table alone, ranked designs before training. Training only on within-site comparisons kept the reverse-set AUC above chance, and when every site was bridged, the test AUC did not drop.

The framework concerns only the study design and the training objective and does not depend on a particular network architecture, so it can be planned before data collection. For data already collected without bridges, it can offer only a site-level evaluation and an explicit statement of the exchangeability assumptions behind any diagnostic claim.

For researchers who pool clinical data, the easiest steps are to tabulate subjects by site and class and to report site-level performance. When collection is still being planned, a few subjects of the missing class can be budgeted at each site and spread evenly; separate case and control collections, if unavoidable, should share scanner and protocol. In either case, a reverse-set stress test should accompany test accuracy. For spurious correlations without bias-conflicting samples in pattern recognition, our results point to a remedy at the level of the data: collecting a few such samples by design. Future work should recruit bridging subjects prospectively at class-pure clinical sites, extend the theory to site-by-diagnosis interactions and non-linear scores, and test within-site objectives for EEG and end-to-end deep models.
